\textbf{(i)} $(\varphi^*)^{-1}(D(f))$ is the set of ideals $\mathfrak q$ such that $\varphi^{-1}(\mathfrak q)$ does not contain $f$. This happens iff $\mathfrak q$ does not contain $\varphi(f)$. So $(\varphi^*)^{-1}(D(f))=D(\varphi(f))$.

\textbf{(ii)} $(\varphi^*)^{-1}(V(\mathfrak a))$ is the set of ideals $\mathfrak q$ such that $\varphi^{-1}(\mathfrak q)$ contains $\mathfrak a$. This is the set of ideals that contains $\varphi(\mathfrak a)$, which are those that contain $V(\mathfrak a^e)$. So $(\varphi^*)^{-1}(V(\mathfrak a))=V(\mathfrak a^e)$.

\textbf{(iii)} Trivially $\varphi^*(V(\mathfrak b))\subseteq V(\mathfrak b^c)$.
\[
\bigcap_{\mathfrak p\in\varphi^*(V(\mathfrak b))}\mathfrak p
=
\bigcap_{\mathfrak q\in V(\mathfrak b)}\varphi^{-1}(\mathfrak q)
=
f^{-1}(\sqrt{\mathfrak b})
=
\sqrt{\mathfrak b^c}
\]
where last equality is by 1.18. This means $\overline{\varphi^*(V(\mathfrak b))}=V(\mathfrak a)$ for some $\mathfrak a\subseteq\sqrt{\mathfrak b^c}$ and implies the other inclusion.

\textbf{(iv)} Since surjective, $\varphi^*$ gives a one to one correspondence between $Y$ and $V(\ker(\varphi))$ by prop 1.1. In particular, it sends closed sets to closed sets, so homeomorphism.

\textbf{(v)} If injective then by (iii), $\overline{\varphi^*(Y)}=V(\varphi^{-1}(0))=V(0)=X$. In general, it is dense iff
\[
\overline{\varphi^*(Y)}=V(\varphi^{-1}(0))=V(\ker\varphi)=X
\]
which happens iff $\ker\varphi$ is contained in nilradical.

\textbf{(vi)} Follows from $(\psi\circ\varphi)^{-1}=\varphi^{-1}\circ\psi^{-1}$.

\textbf{(vii)} $B$ has two prime ideals $(0,K)$ and $(A/\mathfrak p,0)$. The inverse image of them are $\mathfrak p$ and $0$. Thus $\varphi^*$ is bijective. It is not homoemorphism because the singletons in $\operatorname{Spec}B$ are closed, but $\{(0)\}$ is not closed in $\operatorname{Spec}A$.
